\documentclass[12pt]{article}

\usepackage{geometry}
\geometry{
	a4paper,
	margin=2cm,
	left=3cm,
	top=1cm}
\usepackage[utf8]{inputenc}
\usepackage[german]{babel}
\usepackage[T1]{fontenc}
\usepackage{graphicx}
\usepackage{hyperref}

\setlength{\parindent}{0em} %setzt Paragraphen-Einzug auf 0

\title{Angebot "`Interaktives TicTacToe für Kinder"'}

\begin{document}

\begin{center}
\begin{figure}
\includegraphics[scale=0.9]{banner.pdf}
\end{figure}
\end{center}

\small
\begin{flushright}
Kleinprogramm GmbH\\
Eichenweg 6\\
01069 Dresden\\ 
E-Mail: kontakt@kp.de
\end{flushright}

%\vspace*{2mm}

\vspace*{10mm}

\large \textbf{Admindokumentation: \ "`Interaktives\ TicTacToe\ für\ Kinder"'}

\normalsize
\begin{flushright}
Dresden, den 03.06.2019\\
\end{flushright}

\vspace*{7mm}

\tableofcontents
\newpage

\subsection{Installation & Upgrading}
Die Software wird in Rahmen einer weiteren Unterhaltungsmöglichkeit, während der Wartezeit, in der Kinderarztpraxis von Dr. Lasse Niessen verwendet.
Die ausführbare .jar Datei ist sichtbar für die Endanwender auf dem Desktop zu platzieren.
Minimale Updates, die sich ggf. auf Tests zurückführen lassen sind auf Anfrage einzubetten.
Größere Upgrades sind derzeit nur in Form eines weiteren Auftrags einzurichten.
\subsection{Systemspezifische Instruktionen und Pakete}
\subsubsection{Windows}
Das Betriebssystem ist Windows 10 Home 64bit.
\subsubsection{Linux}
Die Lauffähigkeit auf einem Linux System ist nicht nötig, da dieses vom Kunden nicht vorgesehen ist.
\subsection{Security,Performance und Rollen}
Die Software muss sorgfältig und Kundenfreundlich vorinstalliert sein. Effizienz des Programms in sich ist gegeben. Die Anwendung sollte stabil und ohne Unterbrechungen 18 Stunden an 6 Tagen die Woche laufen. Nutzerrollen sind Kinder im Alter zwischen drei und                     14 Jahren, Eltern und Personal einer Arztpraxis.
\subsection{FAQ}
Bei Rückfragen bitte an die oben verwiesene Email Adresse wenden.
\newpage


\vspace*{\fill}
\includegraphics[scale=0.9]{footer.pdf}
\thispagestyle{empty}
\end{document}